\subsection{\textcolor{MidnightBlue}{Erstellung des Fragebogens für die Nutzerbefragung}}
\label{ch:Validierung:sec:Material:Nutzerbefragung}
Der Fragebogen für die Nutzerbefragung wurde nach dem Vorbild von Jochen Prümper (s. \autoref{ch:SdF:sec:Fragebogen}) erstellt. Das Validierungsverfahren zielt auf die Auswertung der Usability der Dashboard-Lösung im Vergleich zum bestehenden System ab. Da in die ISO-Norm jedoch auch die Auswertung der User-Experience inbegriffen ist, wurde der Fragenkatalog gekürzt. Die Abgrenzung zwischen den Begriffen Usability und User-Experience wird in \autoref{ch:SdF:sec:Usability_vs_User_Experience} dargestellt.

Der Fragenkatalog wurde auf zehn Fragen gekürzt, um den Umfang gering zu halten. Die genauen Begründungen werden in \autoref{Add:Nutzerbefragung_Prümper:Begründung_Filterung} aufgelistet.

Der daraus resultierende Fragebogen ist in \autoref{Add:Nutzerbefragung_Prümper:gekürzter_Fragebogen} aufgeführt und dessen Fragen den Interaktionsprinzipien zugeordnet. Die verwendeten Abkürzungen der Prinzipien sind aus dem Originalfragebogen übernommen und in \autoref{Add:Nutzerbefragung_Prümper:ISO-Abkürzungen} aufgeschlüsselt.

Nachfolgend wird mit der Kurzform 'Frage 1', 'Frage 2' etc. auf den Fragebogen referenziert.

%Auf der Grundlage einer faktorenanalytischen Überprüfung wurde mit dem ISONORM 9241/10-S eine Kurzversion des ISONORM 92410/10 erstellt. Dabei konnten die vorhandenen sieben Skalen bestätigt und ihnen je 3 Items zugeordnet werden. \cite{Pruemper:Fragebogen-S:2006:Paper}

%Der Fragebogen wurde nach dem Vorbild der Nutzerb
%- Vorbild nach Prümper; \cite{Pruemper:Fragebogen-S:2006}\\
%- Begründung warum welche Fragen rausgekürzt wurden, s. \autoref{Add:Nutzerbefragung_Prümper:Begründung_Filterung}\\
%- verwendeter Fragebogen, s. \autoref{Add:Nutzerbefragung_Prümper:gekürzter_Fragebogen}
%-> Aufschlüsselung der Fragen zu den Interaktionsprinzipien der ISO-Norm und die verwendete Reihenfolge\\
%-> nachfolgend wird mit der Kurzform 'Frage 1', 'Frage 2' etc. auf den Fragebogen referenziert
